\section{Requisito del Producto}
\normalsize{ \indent
Para ejecutarse la página web de forma correcta, los docentes
necesitarán una conexión mínima de 1 Mbps, ya que cuenta con
componentes mínimos para la carga de imágenes y formulario
digital. Por otro lado, la Rectora y Coordinación Escolar
requieren una conexión mínima de 2,5 Mbps; por suerte, en el
colegio se cuenta con WiFi mucho superior a tal cantidad. El
procesador del dispositivo de navegación debe ser de 4 MB de
RAM como mínimo y cada dispositivo accederá al sitio
\url{https://www.ilc.edu.ar/}
para entrar al sistema (se tendrá un subdominio dedicado para
el ingreso a legajos, cargos y horarios docentes; incluso se
puede utilizar el mismo para la totalidad del proyecto).
}
\section{Requisito Organizacional}
\normalsize{ \indent
Cuando se obtienen los datos de la declaración jurada, la firma
debe ser presencial; de tal forma, se verifica que la firma es
real. Además, tener mayor inversión en comprar un sistema de
firma digital que sólo se usaría en el colegio tendría muchos
costos; en consecuencia, se prefiere que al firmar sea en una
hoja física con tinta.
}
\section{Requisito Externo}
\normalsize{ \indent
Se deberá conservar la privacidad de los datos personales
envueltos por secreto profesional. Probablemente, el autor de
este proyecto deba firmar un contrato de confidencialidad para
tener evidencia sobre la manipulación de información de los
docentes.
}
